\section{Conclusion and Future Work}\label{sec:conclusion}

\gofer{} v4 turns a shell-script training loop into a tested, resumable pipeline shaped
by the \katago{} recipe~\citep{wu2019accelerating}:
\begin{itemize}
  \item \textbf{Data.} Self-play writes compact shards with correctly framed labels. The
        learner trains on a growing window with game-level validation, symmetry
        augmentation, masked multi-head losses and EMA weights.
  \item \textbf{Architectures.} The learner supports a size-agnostic global-pooling network,
        and it can move to a new architecture by distilling from the champion.
  \item \textbf{Loop.} Candidates pass a parity-checked export and a sequential gate, and
        every champion is kept.
\end{itemize}
Measured against v3, the gains we can defend today are correctness, data volume,
robustness, and control of false promotions at the gate (which costs longer gates for
near-neutral candidates). Measured against \katago{}, v4 is a partial adoption with
different engineering priorities: it is built for a project with no compute budget.

\paragraph{Experiment protocol.}
The open questions are empirical, and the pipeline was built to answer them cheaply. We
propose a protocol modelled on \citeauthor{wu2019accelerating}'s ablations: a main run,
plus runs that each change one component, compared by Elo on the pipeline's champion
ladder against cumulative self-play compute (network queries $\times$ blocks $\times$
channels$^2$, as in \citet{wu2019accelerating}). In priority order:
\begin{enumerate}
  \item fast-search rows kept (masked policy) versus discarded;
  \item \code{gpool}-6$\times$64 versus legacy-6$\times$64 at equal self-play compute;
  \item a \code{gpool} student distilled from the v3 champion versus the same network
        trained from scratch, measured in cycles to reach the champion's arena strength;
  \item per-cycle fine-tuning versus continuous training;
  \item EMA on versus off, and recency decay $\lambda\in\{0,2\}$.
\end{enumerate}
Each comparison should be repeated over several seeds. The pipeline's cost estimator
projects the price of each run from measured stage timings, so runs can be sized to the
available budget before any GPU is rented.

\paragraph{Future engineering.}
Following \katago{}, the next steps are:
\begin{itemize}
  \item \katago{}'s PUCT-based policy target pruning, replacing the visit threshold, and
        noise-free fast searches;
  \item a discretised score-belief head;
  \item liberty and ladder input features (a feature-schema version bump);
  \item board-size and komi randomisation in self-play, which the \code{gpool} trunk
        already accommodates.
\end{itemize}
